\begin{statement}[C.4.5; Proved: complete numerical verification through order eight]
All $13$ smooth exact-order roots for $2\le n\le8$ were evaluated
at $110$ decimal digits independently of the algebraic norm formula.
The discrepancy in $N^{12}$ is below $10^{-90}$ at every root;
the period coordinates also verify $nP=O$ to that tolerance.
There are $0,0,1,2,2,4,4$ roots for $n=2,\ldots,8$.
\end{statement}
\begin{proof}
The script \path{Cspecial_torsion_numerics.py} computes the
roots of $F_n$ to $115$ digits and constructs periods of
$y^2=x^3+Ax+B$ using its three roots $e_1,e_2,e_3$. If
$k=(e_2-e_3)/(e_1-e_3)$ and $K$ denotes the complete elliptic
integral with parameter $k$, it uses
\[
 \tau=iK(1-k)/K(k),\qquad
 v=\frac{F\bigl(\arcsin\sqrt{(e_1-e_3)/(-e_3)},k\bigr)}{2K(k)}.
\]
The identity
$\wp(u)=e_3+(e_1-e_3)/\operatorname{sn}^2
(\sqrt{e_1-e_3}\,u,k)$ verifies that this is a zero of $\wp$.
Permutations of the roots are tried until $\operatorname{Im}\tau>0$;
modular reduction and lattice translations improve numerical
conditioning and leave $N$ unchanged. The program then evaluates
the theta quotient directly. It checks both the norm formula and
that the coefficients of $nv$ in the basis $(\tau,1)$ are integral.
The full output gives at least $70$ digits for every $t,\tau,v,N$
and $N^{12}$, including both nonreal conjugates. The largest measured
norm discrepancy is below $2\cdot10^{-108}$. These are numerical
checks of the exact theorem, not numerical proofs of it.
\end{proof}
